\section{Regulación europea y soberanía: separar el momento de la clase de las reglas vigentes}

Cuando un estudiante preguntó por la EU AI Act, la respuesta de Lample reflejó el estado de ese momento: los requisitos generales de transparencia ya estaban a la vista, pero seguía en discusión qué había que divulgar en concreto y cómo se aplicarían las plantillas técnicas. Ese juicio no puede presentarse sin cambios como la situación de 2026. Después, la Comisión Europea precisó que las obligaciones de transparencia y de derechos de autor de los proveedores de modelos general-purpose AI (GPAI) se aplican desde el 2 de agosto de 2025; los modelos introducidos en el mercado antes de esa fecha tienen un periodo de transición hasta el 2 de agosto de 2027, y los modelos que alcanzan el umbral de riesgo sistémico tienen además obligaciones de seguridad adicionales.

\begin{warningbox}{Nota sobre fechas regulatorias; no constituye asesoría legal}
Esta sección solo explica las implicaciones de infraestructura tratadas en el curso. Para saber si un modelo concreto se considera GPAI, quién es el provider, en qué momento se considera placed on the market y cómo se aplican las excepciones de código abierto o las obligaciones por riesgo sistémico, hay que consultar las directrices vigentes de la Comisión Europea y asesoría legal profesional. Las reglas y sus interpretaciones seguirán actualizándose.
\end{warningbox}

\subsection{Por qué la privacy y la reliability cambian los límites del mercado}

A los clientes de finanzas, salud, defensa y seguros no solo les importa la puntuación del modelo: también les importa si los datos salen de su perímetro, si las dependencias del servicio son controlables, si las versiones pueden congelarse y si es posible revertir ante una falla. Aunque la calidad de ingeniería de una API externa sea mayor, una organización puede optar por un private deployment por los límites de responsabilidad y la continuidad del negocio. Al leer la figura, conviene tratar privacy, reliability y control como tres ejes independientes: a una organización puede faltarle solo uno de ellos, o puede estar dispuesta a asumir mayores costos de operación para obtener garantías más fuertes en los tres a la vez.

\lecturefigure{15-privacy-reliability.png}{Lo que compran las organizaciones sensibles es la combinación de privacy, reliability y control.}{Subtítulos locales 35:40--37:10; concepto redibujado.}

\begin{knowledgebox}{Lectura de la figura: la soberanía no es «construirlo todo por cuenta propia»}
La soberanía se acerca más a la sustituibilidad y al poder de decisión: poder elegir el modelo, migrar los datos, auditar las dependencias, congelar versiones y mantener el servicio si el proveedor se retira. Una organización puede usar modelos externos o recursos de código abierto y, al mismo tiempo, reducir la dependencia de un proveedor mediante interfaces estándar, exportación de datos, evaluación de modelos y estrategias con varios proveedores.
\end{knowledgebox}

\teachervoice{En la clase, la motivación europea abarcaba a la vez el talento, la industria y la soberanía. Lample también aclaró que no es experto en regulación y limitó muchos de sus juicios a «too early to say». Estas notas conservan esa salvedad para no presentar la opinión de un emprendedor como si fuera un hecho de política pública.}

\subsection{Resumen de la sección}

La EU AI Act debe entenderse por fechas: en el momento de la clase, los detalles técnicos aún eran inciertos; después, las GPAI obligations entraron en fase de aplicación. Para el diseño de sistemas, la regulación y la soberanía terminan traduciéndose en lineage de datos, documentación, evaluación, controles de seguridad y despliegues auditables.
